\subsection{Proofs} \label{apx:proof}

\subsubsection{Proof of Theorem~\ref{th:l1}}\label{apx:pfTheorem1}
We describe the simulator $\Sim$ in the $\{\F_\oslone, \F_{2PC} \}$-hybrid model
for the case that Party 1 is corrupted.
The simulator and proof of security are analogous in the case that Party 2 is corrupted.

$\Sim$ receives as input
$\vec{w}_1$, the output $i^*$, and $||\vec{w}_1 + \vec{w}_2||_2$.
$\Sim$ samples $r^*$ from a geometric distribution with success probability $p = ||\vec{w}_1 + \vec{w}_2||_2^2$.

$\Sim$ invokes Party 1 on input $\vec{w}_1$.
For $i \in [r^*-1]$, 
Party 1 sends its input to the first three
invocations of 
$\F_\oslone$ and $\Sim$
returns to it three random values in $\mathbb{Z}_n$.
Party 1 sends its input to the second three
invocations of 
$\F_\oslone$ and $\Sim$
returns to it three random values in $\mathbb{Z}_n$.
Party 1 sends its input to the 
$\F_{2PC}$ functionality and $\Sim$
returns to it $\bot$.
For $i = r^*$,
Party 1 sends its input to the first three
invocation of 
$\F_\oslone$ and $\Sim$
returns to it three random values in $\mathbb{Z}_n$.
Party 1 sends its input to the second three
invocations of 
$\F_\oslone$ and $\Sim$
returns to it three random values in $\mathbb{Z}_n$.
Party 1 sends its input to the 
$\F_{2PC}$ functionality and $\Sim$
returns to it $i^*$.

It is clear that the view of Party 1
is identical in the ideal and real world,
assuming that $\Sim$
samples the first succeeding round,
$r^*$, from the correct distribution.
In the following, we argue that this is 
indeed the case.

As was shown in the correctness analysis,
if the protocol has not already halted before
round $r$, then
the probability of halting (and outputting some valid
index) in round $r$ is:
\[
||\vec{w}_1||^2 + 2 \langle \vec{w}_1, \vec{w}_2 \rangle +  ||\vec{w}_2||^2 = ||\vec{w}_1 | \vec{w}_2||_2^2.
\]
Since $r^*$ is defined as the round in which the protocol halts,
the distribution on $r^*$ is exactly the distribution on
the number of Bernoulli trials
(with probability $p = ||\vec{w}_1 | \vec{w}_2||_2^2$) needed to
get one success.  Sampling the number of rounds is therefore equivalent
to sampling the random variable corresponding to the number of rounds
from a geometric distribution with success probability $p = ||\vec{w}_1 | \vec{w}_2||_2^2$,
which is exactly what $\Sim$ does.


\subsubsection{Proof of Theorem~\ref{thm:os-fhe}}\label{apx:pfTheorem2}
We first give the simulation of the sender. The simulator proceeds as
follows: 
\BI
\item Send  $\vec{w}$ to $\F_{\oslone}$ and receive $\pi$ as the output.  Place $\pi$ on the random tape of the sender. 


\item Simulate the output of $F_{2PC}$ by sending a random
$r_1$  to the sender.  
\item Simulate the key generation protocol honestly. 
\item Send a random encryption for
$\lbra r_2 \lket$ in Step 4. 
\EI
Due the semantic security of the underlying FHE scheme, the simulation
is indistinguishable. 
%
Next, we give the simulation of the receiver. 

\BI
\item The simulator receives output $i \xor \pi$ from $\F_{\oslone}$.

\item The simulator works as functionality $F_{2PC}$ sending a random
$r_2$ as the output to the receiver. 

\item The simulator runs the key generation protocol honestly, and stores the threshold decryption key of the sender. 

\item In step 6, the simulator computes $c := \lbra i \xor \pi \lket$ and sends it
to the receiver. 

\item The simulator runs the threshold decryption protocol honestly. 
\EI
The simulation is perfect. 


\subsubsection{Proof of Theorem~\ref{th:l2}} \label{apx:pfTheorem3}
We describe the simulator $\Sim$ in the $\{\F_{L_1}^{ss}, \F_{2PC}\}$-hybrid model
for the case that Party 1 is corrupted.
The simulator and proof of security are analogous in the case that Party 2 is corrupted.

$\Sim$ receives as input
$\vec{w}_1$ and the output $i^*$.
$\Sim$ invokes Party 1 on input $\vec{w}_1$.
For $j \in [B]$, the simulator works as follows:
\BI    
    
\item Upon Party 1 sending its input to $\F_{L_1}$, $\Sim$ returns a uniformly random share $r$\anote{I don't want to specify the space here since this depends on the precision.}

\item In place of the encryption of $w_{2,i_j}$ from Party 2, $\Sim$ sends Party 1 an FHE ciphertext encrypting 0.

\item Upon Party 1 sending its input to $\F_{2PC}$, $\Sim$ returns to Party 1 an FHE ciphertext encrypting 0.
\EI

% Party 1 sends its input to the
% invocations of 
% $\F_{L_1}$ and $\Sim$
% returns to it a random values in.
% $\Sim$ sends to Party 1 an FHE ciphertext
% encrypting $0$, in place of the encryption of 
% $w_{2,i_j}$ from Party 2.
% Party 1 sends its input to $\F_\bcoin$
% and $\Sim$ sends to Party 1 an FHE ciphertext
% encrypting $0$, on behalf of the ideal functionality.
% After the loop completes,
% Party 1 sends its input to $\F_{2PC}$
% and $\Sim$ returns to it $i^*$
% on behalf of the ideal functionality.

The only differences in the view of Party 1
in the ideal and hybrid worlds,
are that 
(1) In the hybrid world it gets a secret share of
$i_j$, whereas in the ideal world it gets a uniformly random value;
(2) In the hybrid world it gets an encryption of $w_{2,i_j}$ from Party 2, whereas in the ideal world it gets an encryption of 0;
(3) In the hybrid world it gets encryptions of
$i_j$ or $0$ from the ideal functionality
$\F_{2PC}$, whereas in the ideal world it always
gets encryptions of $0$.

Receiving uniformly random values instead of correct secret shares does not affect the view of Party 1, since the additive secret sharing used
has perfect secrecy.
Further, switching from encryptions
of $w_{2,i_j}$ and $i_j$ to encryptions of $0$
is indistinguishable due to the semantic security
of the threshold FHE scheme.
Thus, the view of Party 1 is computationally indistinguishable in the hybrid world and the ideal world.

This concludes the proof of security of the $L_2$
sampling protocol.
 
%
%\subsubsection{Proof of Theorem~\ref{th:l2}}
%Note that the entire view of Party $P_b$ consists of the secret shares
%received in from the oblivious sampling functionality, its view in the
%secure protocol computation, the final output, denoted $i^*$, and the
%round $r^*$ which returned $i^*$.
%
%The simulator receives $i^*$ from the ideal functionality and must draw
%$i^*$ from the correct distribution.  Given this, the entire view of
%Party $P_b$ can be simulated.
%
%As before, we show that the simulator can sample the first succeeding
%round $r^*$ from the exact correct distribution, given only the value of
%$p = ||\vec{w}_1||^2 + 2 \langle \vec{w}_1, \vec{w}_2 \rangle +
%||\vec{w}_2||^2$.
%
%Let ${\Success}$ denote an event in which the protocol succeeds
%in the current round. Note the protocol is engineered to have the
%following property: $\Pr[\Success] = p/6$. 
%Let $\Output(i)$ denote an event in which the protocol outputs
%Then, we have $$q_{i^*} = \Pr[ \Output(i^*) |\Success] =
%\frac{(w_{1,i^*} + w_{2,i^*})^2}{p}.$$
%
%Let $\FirstSuccess(r)$ denote an event in which the protocol
%succeeds for the first time on the $r$-th round.  Then, note that for $r
%\in \mathbb{N}$, the probability that the protocol succeeds for the
%first time in the $r$-th round AND the output
%being $i^*$ is equal to:
%\begin{align*}
%&\Pr[\FirstSuccess(r^*) \wedge \Output(i^*)]\\
%&=
%\Pr[\bot \mbox{ returned in first } r^*-1 \mbox{ rounds}] \cdot
%\Pr[\Success \wedge \Output(i^*)]\\
%&= (1-p/6)^{r^*-1} \cdot p/6 \cdot q_{i^*},
%\end{align*}
%
%Now, the probability that the protocol eventually outputs $i^*$ (after some number of rounds) is:
%\[
%\Pr[\Output(i^*)]
%= \sum_{r=1}^\infty \Pr[\FirstSuccess(r) \wedge \Output(i^*)]
%= q_{i^*}.
%\]
%
%Thus, the \emph{conditional probability} that the protocol succeeds for the first time on the $r^*$-th round conditioned on the output
%being $i^*$ is equal to
%\begin{align*}
%&\Pr[\FirstSuccess(r^*) ~|~ \Output(i^*)] \\
%&= \frac{\Pr[\FirstSuccess(r^*) \wedge \Output(i^*)]}{\Pr[\Output(i^*)] }\\
%&= (1-p/6)^{r^*-1} \cdot p/6.
%\end{align*}
%
%The above is exactly the probability of the number of Bernoulli trials (with probability $p/6$ needed to get one success.
%Sampling the number of rounds is therefore equivalent to sampling the random variable
%corresponding to the number of rounds from
%a geometric distribution with success probability $p/6$.
%This can be efficiently simulated given just the $L_2$ norm $p = ||\vec{w}_1||^2 + 2 \langle \vec{w}_1, \vec{w}_2 \rangle +  ||\vec{w}_2||^2$.
%

\subsubsection{Proof of Lemma~\ref{lem:close_dist_2}}\label{apx:pfLemma5}
We want to show that
\begin{align*}
    \Pr_{D_{L_p}(\vec{w}_1, \vec{w}_2)}[i] &= \frac{(w_{1,i} + w_{2,i})^p}{||\vec{w}_1 + \vec{w}_2||_p^p}\\
   &= \frac{(w_{1,i} + w_{2,i})^p}{c \cdot (||\vec{w}_1||_p^p  + ||\vec{w}_2||_p^p)}\\
    &\leq \frac{2^{p-1} (w^p_{1,i} + w^p_{2,i})}{c \cdot (||\vec{w}_1||_p^p  + ||\vec{w}_2||_p^p)}\\
    &= 2^{p-1}/c \cdot \Pr_{D_{\mathsf{ignore}, p}(\vec{w}_1, \vec{w}_2)}[i].
\end{align*}
The inequality holds due to Jensen's inequality with convex function $f(x) = x^p$:
\begin{align*}
1/2^p \cdot(w_{1,i} + w_{2,i})^p & = f(1/2 \cdot w_{i,1} + 1/2 \cdot w_{i,2}) \\
& \le 1/2 \cdot f(w_{i,1}) +  1/2 \cdot f(w_{2,i}) \\
& = 1/2 (w_{i,1}^p + w_{i,2}^p). 
\end{align*}
This completes the proof of Lemma~\ref{lem:close_dist_2}.
 

\subsubsection{Proof of Lemma~\ref{lem:lp}}\label{apx:pfLemma6}
Note that $\Pi_{L_p}$ simply performs rejection sampling in a distributed setting where sampling from $D_{\mathsf{ignore}, p}(\vec{w}_1, \vec{w}_2)$ and computing the probabilities is done in a distributed manner. It is therefore well-known that
as long as for all $i \in [n]$,
\begin{equation} \label{eq:rejection_2}
    \Pr_{D_{L_p}(\vec{w}_1, \vec{w}_2)}[i] \leq 2^p/c \cdot \Pr_{D_{\mathsf{ignore}, p}(\vec{w}_1, \vec{w}_2)}[i],
\end{equation}
then $\Pi_{L_2}$ samples from the exact correct distribution, and
the number of samples required from $D_{\mathsf{ignore}, p}(\vec{w}_1, \vec{w}_2)$ in protocol $\Pi_{L_2}$ follows a geometric distribution with probability $c/2^p$.
Thus, if condition (\ref{eq:rejection_2}) is met, the protocol 
samples exactly correctly and completes in an expected $2^p/c \leq 2^p \in \tilde{O}(1)$ number of rounds (since $c \geq 1$ and $p \in O(1)$).
Further, it can be immediately noted that condition (\ref{eq:rejection_2}) is met due to Lemma~\ref{lem:close_dist_2}.
Finally, each round has $\tilde{O}(1)$ communication, since $\Pi_{\mathsf{ignore}}$ has communication $\tilde{O}(1)$ (by Lemma~\ref{lem:ignore_2})
and since, in addition to that, only a constant number of length $\tilde{O}(1)$ values are exchanged in each round.
Combining the above, we have that $\Pi_{L_p}$ has expected
communication $\tilde{O}(1)$
and worst case (with all but negligible probability) communication $\tilde{O}(1)$.

\subsubsection{Proof of Lemma~\ref{lem:simip}}\label{apx:pfLemma10}
We will show that the joint distribution over the $i$-th entries of
$\tilde{\vec{w}_1} := \vec{M} \w_1 = (\tilde w_{1,1}, \ldots, \tilde w_{1,k})$,
$\tilde{\vec{w}_2} = \vec{M} \w_2 = (\tilde w_{2,1}, \ldots, \tilde w_{2,k})$
can be sampled perfectly, given only $\langle \vec{w}_1, \vec{w}_1 \rangle$,
$\langle \vec{w}_2, \vec{w}_2 \rangle$, and $\langle \vec{w}_1, \vec{w}_2
\rangle$.  

Due to independence of each of the coordinates of $\tilde{\vec{w}_1},
\tilde{\vec{w}_2}$, this immediately implies that the entire
$\mathsf{approxIP}(\vec{w}_1, \vec{w}_2)$ can be simulated perfectly given only
$\langle \vec{w}_1, \vec{w}_1 \rangle$, $\langle \vec{w}_2, \vec{w}_2 \rangle$,
and $\langle \vec{w}_1, \vec{w}_2 \rangle$.

\noindent
We begin by noting that
\begin{eqnarray*}
\tilde{w}_{1,i}  &=& w_{1,1} M_{i,1} + w_{1,2} M_{i,2} + \cdots + w_{1,n} M_{i,n} \\
\tilde{w}_{2,i}  &=& w_{2,1} M_{i,1} + w_{2,2} M_{i,2} + \cdot + w_{2,n} M_{i,n}
\end{eqnarray*}

\noindent
In the following, we show how to jointly sample
$(\tilde{w}_{1,i}, \tilde{w}_{2,i})$.

\begin{description}
\item{Step 1:}
We begin by sampling from the marginal distribution over the first element of
the tuple $\tilde{w}_{1,i}$.  Note that $\tilde{w}_{1,i}$ is distributed
exactly as a Gaussian random variable with mean $0$ and variance $\langle
\vec{w}_1, \vec{w}_1 \rangle$.  Thus, we can perfectly sample from the marginal
distribution over $\tilde{w}_{1,i}$ given only $\langle \vec{w}_1, \vec{w}_1
\rangle$.  Let $z$ be the resulting sample.

\item{Step 2:} We would now like to sample from the conditional distribution
$\tilde{w}_{2,i}$, conditioned on $\tilde{w}_{1,i} = z$.

First, the conditional distribution of $M_{i,1}, \ldots, M_{i,n}$ conditioned
on $w_{1,1} M_{i,1} + w_{1,2} M_{i,2} + \cdots + w_{1,n} M_{i,n} = z$ is
defined by the multivariate Gaussian distribution with the following mean
$\vec{\mu}$ and covariance matrix $\vec{\Sigma}$
(see Corollary 7 in~\cite{C:DDGR20}):
\[
\vec{\mu} = \frac{z}{\langle \vec{w}_1, \vec{w}_1 \rangle} \cdot \vec{w}_1
, ~~~~~
\vec{\Sigma} = \vec{I} - \frac{\vec{w}_1\vec{w}_1^T}{\langle \vec{w}_1, \vec{w}_1 \rangle}
\]
Now, $\tilde{w}_{2,i}$ is a linear combination of the variables $M_{i,1},
\ldots, M_{i,n}$ with coefficients $w_{2,1}, \ldots, w_{2,n}$.  Therefore,
$\tilde{w}_{2,i}$ is distributed as a univariate Gaussian with mean $\mu'$ and
variance $\sigma'$ as follows (see \cite{Liu} for example).
\[
\mu' = \langle \vec{w}_2, \vec{\mu} \rangle = \frac{z \langle \vec{w}_1,
\vec{w}_2 \rangle}{\langle \vec{w}_1, \vec{w}_1 \rangle} 
\]
and
\[
\sigma' = \vec{w}_2^T\vec{\Sigma}\vec{w}_2
= \vec{w}_2^T\vec{w}_2 - \frac{\vec{w}_2^T\vec{w}_1\vec{w}_1^T\vec{w}_2}{\langle \vec{w}_1, \vec{w}_1 \rangle}
= \langle \vec{w}_2, \vec{w}_2 \rangle -
\frac{(\langle \vec{w}_1, \vec{w}_2 \rangle)^2}{\langle \vec{w}_1, \vec{w}_1 \rangle}
\]
Note that the mean and variance depend only on 
$\langle \vec{w}_1, \vec{w}_1 \rangle$, 
$\langle \vec{w}_2, \vec{w}_2 \rangle$,
and $\langle \vec{w}_1, \vec{w}_2 \rangle$,
so we can sample from this distribution given only
those values.
Let $y$ be the resulting sample.
\item{Step 3:}
Output $(z, y)$
\end{description}


 